\documentclass[10pt,a4paper]{amsart}
\usepackage{amsmath,amssymb,amsthm,graphicx}
% Retained preamble for the Stacks Project chapter snapshot in this directory.
% The chapter's own preamble.tex is a copy of the shared Stacks Project preamble; the block
% below is sliced MECHANICALLY, line for line, from it. Nothing was added, removed or
% redefined; the wrapper class, the xr-hyper/hyperref wiring and the externaldocument list
% are omitted because the assembler supplies the document class.
%
% lines taken from preamble.tex: 3, 5, 9, 11, 12, 156, 157, 158, ... (51 lines total)
%
% package coverage restored from the chapter's own preamble (2026-10-13): xy, xy-2cell

\usepackage[all]{xy}
\xyoption{2cell}
\UseAllTwocells

\newenvironment{reference}{\comment}{\endcomment}
\newenvironment{slogan}{\comment}{\endcomment}
\newenvironment{history}{\comment}{\endcomment}
\theoremstyle{plain}
\newtheorem{theorem}[subsection]{Theorem}
\newtheorem{proposition}[subsection]{Proposition}
\newtheorem{lemma}[subsection]{Lemma}
\theoremstyle{definition}
\newtheorem{definition}[subsection]{Definition}
\newtheorem{example}[subsection]{Example}
\newtheorem{exercise}[subsection]{Exercise}
\newtheorem{situation}[subsection]{Situation}
\theoremstyle{remark}
\newtheorem{remark}[subsection]{Remark}
\newtheorem{remarks}[subsection]{Remarks}
\numberwithin{equation}{subsection}
\def\lim{\mathop{\mathrm{lim}}\nolimits}
\def\colim{\mathop{\mathrm{colim}}\nolimits}
\def\Spec{\mathop{\mathrm{Spec}}}
\def\Hom{\mathop{\mathrm{Hom}}\nolimits}
\def\Ext{\mathop{\mathrm{Ext}}\nolimits}
\def\SheafHom{\mathop{\mathcal{H}\!\mathit{om}}\nolimits}
\def\SheafExt{\mathop{\mathcal{E}\!\mathit{xt}}\nolimits}
\def\Sch{\mathit{Sch}}
\def\Mor{\mathop{\mathrm{Mor}}\nolimits}
\def\Ob{\mathop{\mathrm{Ob}}\nolimits}
\def\Sh{\mathop{\mathit{Sh}}\nolimits}
\def\NL{\mathop{N\!L}\nolimits}
\def\CH{\mathop{\mathrm{CH}}\nolimits}
\def\proetale{{pro\text{-}\acute{e}tale}}
\def\etale{{\acute{e}tale}}
\def\QCoh{\mathit{QCoh}}
\def\Ker{\mathop{\mathrm{Ker}}}
\def\Im{\mathop{\mathrm{Im}}}
\def\Coker{\mathop{\mathrm{Coker}}}
\def\Coim{\mathop{\mathrm{Coim}}}
\DeclareMathSymbol{\boxtimes}{\mathbin}{AMSa}{"02}
\def\QCohstack{\mathcal{QC}\!\mathit{oh}}
\def\Cohstack{\mathcal{C}\!\mathit{oh}}
\def\Spacesstack{\mathcal{S}\!\mathit{paces}}
\def\Quotfunctor{\mathrm{Quot}}
\def\Hilbfunctor{\mathrm{Hilb}}
\def\Curvesstack{\mathcal{C}\!\mathit{urves}}
\def\Polarizedstack{\mathcal{P}\!\mathit{olarized}}
\def\Complexesstack{\mathcal{C}\!\mathit{omplexes}}
\def\Pic{\mathop{\mathrm{Pic}}\nolimits}
\def\Picardstack{\mathcal{P}\!\mathit{ic}}
\def\Picardfunctor{\mathrm{Pic}}
\def\Deformationcategory{\mathcal{D}\!\mathit{ef}}
\usepackage{cleveref}
\crefname{originallemma}{original lemma}{original lemmas}
\Crefname{originallemma}{Original lemma}{Original lemmas}
\makeatletter
\newlabel{duality-lemma-affine-flat-Noetherian-gorenstein}{{0C08}{source}{}{sourcepointer-1}{}}
\newlabel{duality-lemma-affine-flat-Noetherian-gorenstein@cref}{{[originallemma][0][]0C08}{[1][0][]source}{}{}{}}
\makeatother
\crefname{originallemma}{original lemma}{original lemmas}
\Crefname{originallemma}{Original lemma}{Original lemmas}
\makeatletter
\newlabel{duality-lemma-flat-quasi-finite-shriek}{{0BV7}{source}{}{sourcepointer-2}{}}
\newlabel{duality-lemma-flat-quasi-finite-shriek@cref}{{[originallemma][0][]0BV7}{[1][0][]source}{}{}{}}
\makeatother
\crefname{originallemma}{original lemma}{original lemmas}
\Crefname{originallemma}{Original lemma}{Original lemmas}
\makeatletter
\newlabel{dualizing-lemma-dualizing-finite}{{0AX0}{source}{}{sourcepointer-3}{}}
\newlabel{dualizing-lemma-dualizing-finite@cref}{{[originallemma][0][]0AX0}{[1][0][]source}{}{}{}}
\makeatother
\crefname{originallemma}{original lemma}{original lemmas}
\Crefname{originallemma}{Original lemma}{Original lemmas}
\makeatletter
\newlabel{dualizing-lemma-gorenstein}{{0DW9}{source}{}{sourcepointer-4}{}}
\newlabel{dualizing-lemma-gorenstein@cref}{{[originallemma][0][]0DW9}{[1][0][]source}{}{}{}}
\makeatother
\crefname{originallemma}{original lemma}{original lemmas}
\Crefname{originallemma}{Original lemma}{Original lemmas}
\makeatletter
\newlabel{lemma-characterize-invertible}{{0DWK}{source}{}{sourcepointer-5}{}}
\newlabel{lemma-characterize-invertible@cref}{{[originallemma][0][]0DWK}{[1][0][]source}{}{}{}}
\makeatother
\crefname{originallemma}{original lemma}{original lemmas}
\Crefname{originallemma}{Original lemma}{Original lemmas}
\makeatletter
\newlabel{lemma-compare-dualizing}{{0BUL}{source}{}{sourcepointer-6}{}}
\newlabel{lemma-compare-dualizing@cref}{{[originallemma][0][]0BUL}{[1][0][]source}{}{}{}}
\makeatother
\crefname{originallemma}{original lemma}{original lemmas}
\Crefname{originallemma}{Original lemma}{Original lemmas}
\makeatletter
\newlabel{remark-relative-dualizing-for-quasi-finite}{{0BVG}{source}{}{sourcepointer-7}{}}
\newlabel{remark-relative-dualizing-for-quasi-finite@cref}{{[originallemma][0][]0BVG}{[1][0][]source}{}{}{}}
\makeatother
\title{Stacks, Discriminants: lemma-gorenstein-quasi-finite}
\author{The Stacks Project Authors}
\date{stacks-project:discriminant.tex:local-snapshot-20261012; GNU FDL 1.2 or later}
\begin{document}
\maketitle
\noindent\textit{Editorial note.} Statement span is lines 3507-3521 of sources/stacks-discriminant/pinned-original-source.tex; proof span is lines 3523-3550, the complete author proof (26 body lines).\par\medskip
\noindent\textit{Editorial note.} Navigation stubs: 7 cited-but-not-included labels are supplied as non-load-bearing pointers with their pinned tags where the snapshot tag index has them.\par\medskip
\noindent\textit{Editorial note.} extra\_packages is empty on purpose: the retained preamble already loads xy with options and a second declaration clashes in the assembler wrapper.\par\medskip
\noindent\textit{Editorial note.} Text is retained verbatim, including the author's own typos and placeholders.\par\medskip
\noindent\textit{Editorial note.} Queue index 21 of queue\_discriminant.json (pb=1285, stmt 3507-3521, proof 3523-3550).\par\medskip
\noindent\textit{Allowed standard prerequisites (editorial).}\begin{itemize}
\item Explicit prerequisite boundary. The cited inputs are prerequisites; the argument in the retained span is not replaced by any cited result.
\end{itemize}\medskip
\ifdefined\hypertarget\hypertarget{sourcepointer-1}{}\fi
\noindent\textit{Source reference.} Original lemma 0C08: sources/stacks-discriminant/pinned-original-source.tex, source lines 3550--3550. Cited by the retained statement or proof but not part of the collected spans; supplied only as a navigation pointer. The argument does not depend on its content.\par\medskip
\ifdefined\hypertarget\hypertarget{sourcepointer-2}{}\fi
\noindent\textit{Source reference.} Original lemma 0BV7: sources/stacks-discriminant/pinned-original-source.tex, source lines 3550--3550. Cited by the retained statement or proof but not part of the collected spans; supplied only as a navigation pointer. The argument does not depend on its content.\par\medskip
\ifdefined\hypertarget\hypertarget{sourcepointer-3}{}\fi
\noindent\textit{Source reference.} Original lemma 0AX0: sources/stacks-discriminant/pinned-original-source.tex, source lines 3550--3550. Cited by the retained statement or proof but not part of the collected spans; supplied only as a navigation pointer. The argument does not depend on its content.\par\medskip
\ifdefined\hypertarget\hypertarget{sourcepointer-4}{}\fi
\noindent\textit{Source reference.} Original lemma 0DW9: sources/stacks-discriminant/pinned-original-source.tex, source lines 3550--3550. Cited by the retained statement or proof but not part of the collected spans; supplied only as a navigation pointer. The argument does not depend on its content.\par\medskip
\ifdefined\hypertarget\hypertarget{sourcepointer-5}{}\fi
\noindent\textit{Source reference.} Original lemma 0DWK: sources/stacks-discriminant/pinned-original-source.tex, source lines 2092--2103. Cited by the retained statement or proof but not part of the collected spans; supplied only as a navigation pointer. The argument does not depend on its content.\par\medskip
\ifdefined\hypertarget\hypertarget{sourcepointer-6}{}\fi
\noindent\textit{Source reference.} Original lemma 0BUL: sources/stacks-discriminant/pinned-original-source.tex, source lines 3360--3375. Cited by the retained statement or proof but not part of the collected spans; supplied only as a navigation pointer. The argument does not depend on its content.\par\medskip
\ifdefined\hypertarget\hypertarget{sourcepointer-7}{}\fi
\noindent\textit{Source reference.} Original lemma 0BVG: sources/stacks-discriminant/pinned-original-source.tex, source lines 649--660. Cited by the retained statement or proof but not part of the collected spans; supplied only as a navigation pointer. The argument does not depend on its content.\par\medskip
\begin{lemma}
\label{lemma-gorenstein-quasi-finite}
Let $f : Y \to X$ be a quasi-finite morphism of Noetherian schemes.
The following are equivalent
\begin{enumerate}
\item $f$ is Gorenstein,
\item $f$ is flat and the fibres of $f$ are Gorenstein,
\item $f$ is flat and $\omega_{Y/X}$ is invertible
(Remark \ref{remark-relative-dualizing-for-quasi-finite}),
\item for every $y \in Y$ there are affine opens
$y \in V = \Spec(B) \subset Y$, $U = \Spec(A) \subset X$
with $f(V) \subset U$ such that $A \to B$ is flat
and $\omega_{B/A}$ is an invertible $B$-module.
\end{enumerate}
\end{lemma}

\begin{proof}
Parts (1) and (2) are equivalent by definition. Parts (3) and (4)
are equivalent by the construction of $\omega_{Y/X}$ in
Remark \ref{remark-relative-dualizing-for-quasi-finite}.
Thus we have to show that (1)-(2) is equivalent to (3)-(4).

\medskip\noindent
First proof. Working affine locally we can assume $f$ is a separated
morphism and apply Lemma \ref{lemma-compare-dualizing} to see that
$\omega_{Y/X}$ is the zeroth cohomology sheaf of $f^!\mathcal{O}_X$.
Under both assumptions $f$ is flat and quasi-finite, hence
$f^!\mathcal{O}_X$ is isomorphic to $\omega_{Y/X}[0]$, see
Duality for Schemes, Lemma \ref{duality-lemma-flat-quasi-finite-shriek}. Hence
the equivalence follows from
Duality for Schemes, Lemma
\ref{duality-lemma-affine-flat-Noetherian-gorenstein}.

\medskip\noindent
Second proof. By Lemma \ref{lemma-characterize-invertible},
we see that it suffices to prove the equivalence of
(2) and (3) when $X$ is the spectrum of a field $k$.
Then $Y = \Spec(B)$ where $B$ is a finite $k$-algebra.
In this case $\omega_{B/A} = \omega_{B/k} = \Hom_k(B, k)$
placed in degree $0$ is a dualizing complex for $B$, see
Dualizing Complexes, Lemma \ref{dualizing-lemma-dualizing-finite}.
Thus the equivalence follows from
Dualizing Complexes, Lemma \ref{dualizing-lemma-gorenstein}.
\end{proof}

\end{document}
